\chapter{Практика}

\section{Сравнение полиномиальных регрессий разных порядков}

Сравним полиномиаьные регрессии с порядком от 2 до 5, чтобы увидеть, как сильно порядок влияет на точность

\begin{figure}[h!]
    \centering

    \begin{subfigure}{0.45\textwidth}
        \centering
        \includegraphics[width=\linewidth]{poly_reg_N=2.png}
        \caption{N = 2}
    \end{subfigure}
    \hfill
    \begin{subfigure}{0.45\textwidth}
        \centering
        \includegraphics[width=\linewidth]{poly_reg_N=3.png}
        \caption{N = 3}
    \end{subfigure}

    \vspace{0.5cm}

    \begin{subfigure}{0.45\textwidth}
        \centering
        \includegraphics[width=\linewidth]{poly_reg_N=4.png}
        \caption{N = 4}
    \end{subfigure}
    \hfill
    \begin{subfigure}{0.45\textwidth}
        \centering
        \includegraphics[width=\linewidth]{poly_reg_N=5.png}
        \caption{N = 5}
    \end{subfigure}

    \caption{Сравнение полиномиальной регрессии при разных степенях}
\end{figure}

\begin{figure}[H]
    \centering
    \includegraphics{acc.png}
    \caption{Точности моделей}
\end{figure}

Точность меняется на уровне погрешности. Это говорит о том, что не для всех случаев порядок полиномиальной регрессии может повышать точность.

\endinput